\documentclass{article}
\usepackage{amsmath, amssymb, amsthm}
\usepackage{hyperref}
\usepackage{geometry}
\geometry{margin=1in}

\title{Erd\H{o}s Problem \#793}
\date{Last updated: 2025-08-31}
\author{Source: erdosproblems.com}

\begin{document}
\maketitle

\noindent\textbf{Status:} OPEN \quad \textbf{No prize}

\noindent\textbf{Tags:} number theory

\medskip
\noindent\textbf{Problem Statement:}

Let $F(n)$ be the maximum possible size of a subset $A\subseteq\{1,\ldots,n\}$ such that $a\nmid bc$ whenever $a,b,c\in A$ with $a\neq b$ and $a\neq c$. Is there a constant $C$ such that\[F(n)=\pi(n)+(C+o(1))n^{2/3}(\log n)^{-2}?\]

\bigskip
\noindent\textbf{Remarks:}

Erd\H{o}s \cite{Er38} proved there exist constants $0<c_1\leq c_2$ such that\[\pi(n)+c_1n^{2/3}(\log n)^{-2}\leq F(n) \leq \pi(n)+c_2n^{2/3}(\log n)^{-2}.\]Erd\H{o}s \cite{Er69} gave a simple proof that $F(n) \leq \pi(n)+n^{2/3}$: define a graph with vertex set the union of those integers in $[1,n^{2/3}]$ with all primes $p\in (n^{2/3},n]$. We have an edge $u\sim v$ if and only if $uv\in A$. It is easy to see that every $m\leq n$ can be written as $uv$ where $u\leq n^{2/3}$ and $v$ is either prime or $\leq n^{2/3}$, and hence there are $\geq \lvert A\rvert$ many edges. This graph contains no path of length $3$ and hence must be a tree and have fewer edges than vertices, and we are done. This can be improved to give the upper bound mentioned by using a subset of integers in $[1,n^{2/3}]$.

More generally, one can ask for such an asymptotic for the size of sets such that no $a\in A$ divides the product of $r$ distinct other elements of $A$, with the exponent $2/3$ replaced by $\frac{2}{r+1}$.

For further discussion and references concerning this generalisation to $r\geq 3$ see the comment by Wouter van Doorn.

See also [425].

\bigskip
\noindent\textbf{References:}

[Er38] P. Erd\H{o}s, On sequences of integers no one of which divides the product of two others and on related problems. Tomsk. Gos. Univ. Ucen Zap. (1938), 74-82.
 
 [Er69] Erd\H{o}s, Paul, Some applications of graph theory to number theory. The Many Facets of Graph Theory (Proc. Conf.,
Western Mich. Univ., Kalamazoo, Mich., 1968) (1969), 77-82.

\bigskip
\noindent\small{Source: \url{https://www.erdosproblems.com/793}}
\end{document}
